% Generated draft: requires mathematical and source review.
\documentclass[11pt]{article}
\usepackage[margin=25mm]{geometry}
\usepackage{fontspec}
\setmainfont{Latin Modern Roman}
\usepackage{amsmath,amssymb,mathtools,unicode-math}
\setmathfont{Latin Modern Math}
\usepackage{enumitem,needspace,xurl,hyperref}
\setlength{\parindent}{0pt}
\setlength{\parskip}{5pt}
\setlength{\emergencystretch}{3em}
\newcommand{\N}{\mathbb{N}}
\newcommand{\Z}{\mathbb{Z}}
\newcommand{\Q}{\mathbb{Q}}
\newcommand{\R}{\mathbb{R}}
\newcommand{\C}{\mathbb{C}}
\newcommand{\F}{\mathbb{F}}
\newcommand{\A}{\mathbb{A}}
\newcommand{\PP}{\mathbb{P}}
\newcommand{\E}{\mathbb{E}}
\newcommand{\eps}{\varepsilon}
\newcommand{\dd}{\,\mathrm{d}}
\begin{document}
\section*{UnsolvedMath ID: 30001017}
\textbf{Source problem code:} catalogue entry
\subsection{Definitions and mathematical statement}
For a positive integer $g$, let $\mathcal A_g$ be the moduli space of principally polarized abelian varieties of dimension $g$, and let $\mathcal A_g^F$ be its first Voronoi, or perfect cone, compactification. Put $N=g(g+1)/2=\dim\mathcal A_g^F$. Let $L$ be the Hodge divisor class, conventionally the first Chern class of the determinant Hodge bundle, and let $D$ be the source's boundary divisor class. With the appropriate rational intersection-theory convention define \[a_n^{(g)}=\deg_{\mathcal A_g^F}\bigl(L^nD^{N-n}\bigr),\qquad n\in\Z,\quad0\leq n\leq N.\] The conjecture is \[n\notin\left\{\frac{k(k+1)}2:k\in\Z,\ 0\leq k\leq g\right\}\quad\Longrightarrow\quad a_n^{(g)}=0.\] Equivalently, a nonzero top intersection number may occur only at a triangular exponent of $L$. This is only a necessary condition for nonvanishing; it does not assert that every listed triangular exponent gives a nonzero value. The compactification is the first Voronoi/perfect cone one, not a different toroidal compactification. Stack versus coarse-space normalization and the boundary multiplicity convention are not specified; their role should be confirmed even though the target here is vanishing. The supplied 2026-08-21 triage labels the assertion open and reports computations in substantial ranges and low genus without a general proof.
\subsection{Short English statement}
Must every top Hodge-boundary intersection vanish unless the exponent of the Hodge class is a triangular number? The question gives a necessary pattern for nonvanishing, without asserting its converse.
\subsection{Sources}
\begin{itemize}\item Local UnsolvedMath registry, permanent ID 30001017, source problem catalogue entry, \emph{Intersection Numbers on First Voronoi Compactifications}, proposed year 2008; no individual proposer supplied.\item Supplied citation: \emph{Komplexe Analysis} (2009), Oberwolfach Reports, DOI 10.4171/owr/2008/38, pp.\ 2165--2218: \url{https://doi.org/10.4171/owr/2008/38}. Supplied literature label: original report.\item Supplied literature label: Erdenberger, Grushevsky and Hulek, \url{https://arxiv.org/abs/0707.1274}.\item Literature assessment date: 2026-08-21. The status and progress claims are supplied dated curation, not independently verified here.\end{itemize}
\end{document}
